\documentclass{article}

\usepackage{iclr2027_conference,times}
\usepackage[T1]{fontenc}
\usepackage{lmodern}
\usepackage{amsmath,amssymb,amsthm,mathtools}
\usepackage{booktabs,microtype,url,xcolor,graphicx}
\usepackage[colorlinks=true,allcolors=blue]{hyperref}

\newtheorem{theorem}{Theorem}
\newtheorem{lemma}[theorem]{Lemma}
\newtheorem{proposition}[theorem]{Proposition}
\newtheorem{corollary}[theorem]{Corollary}
\theoremstyle{remark}
\newtheorem{remark}[theorem]{Remark}

\newcommand{\E}{\mathbb E}
\newcommand{\R}{\mathbb R}
\newcommand{\Risk}{\mathfrak R}
\newcommand{\Stoch}{\mathcal S}
\newcommand{\1}{\mathbf 1}
\newcommand{\RR}{\textsf{RR}}
\newcommand{\SGD}{\textsf{SGD}}

\title{The Missing Horizon in Random Reshuffling}
\author{\textbf{Pahan Dewasurendra} \\ Johns Hopkins University}

\begin{document}
\maketitle

\begin{abstract}
Random reshuffling now has an upper bound that dominates stochastic gradient
descent at every stable constant step size and every finite horizon. Its
matching lower bound was left open. We solve this problem for the exact
all-inner average used by the upper bound. The answer is not the published
horizon-independent stochastic envelope. If
$\alpha=\eta nL\leq1$, the fixed-step minimax stochastic error is
\[
 \Theta\!\left(\eta^2nL\sigma_\star^2
 \left[K^{-1}+\min\{(\alpha K)^2,1\}\right]\right).
\]
It decreases, increases, and then plateaus across three horizon regimes. Its
minimum occurs at $K\asymp\alpha^{-2/3}$, before the usual mixing scale
$\alpha^{-1}$. For $\alpha\geq1$, the law is
$\Theta(\eta\sigma_\star^2)$. The upper bound follows from an exact
finite-population bridge identity and a dimension-free control of all state
memory as a Lipschitz remainder. The lower bound combines a centered
common-Hessian bridge with a rectifying piecewise quadratic in two dimensions.
An exact telescoping identity explains why common-Hessian quadratics alone
miss the eventual floor. Together with the deterministic term, this gives a
complete minimax characterization for all horizons and all stable constant
steps. Optimizing the step recovers the best known tuned rate, while the
fixed-step theorem reveals a nonmonotone transient hidden by that rate.
\end{abstract}

\section{Introduction}

Random reshuffling processes every component of a finite sum exactly once per
epoch in a fresh random order. It is the common practical implementation of
stochastic gradient descent, and it is often faster than sampling with
replacement. A long line of work has sought to explain this advantage
\citep{gurbuzbalaban2021why,nagaraj2019sgd,haochen2019random,
mishchenko2020random,nguyen2021unified,safran2020how,safran2021many}.

The recent result of \citet{liu2026dominates} resolves the upper-bound side for
smooth convex finite sums. For any stable step $\eta\lesssim1/L$ and any
number of epochs $K$, it bounds the exact all-inner average by
\begin{equation}
 \frac{D^2}{\eta nK}
 +\min\{\eta,\eta^2nL\}\sigma_\star^2 .       \label{eq:liu-envelope}
\end{equation}
This is never slower than the corresponding bound for with-replacement
\SGD. The same paper identifies the matching lower characterization for all
steps and horizons as an important open problem. Existing lower bounds cover
small steps, sufficiently many epochs, and averages of epoch-boundary
iterates \citep{cha2023tighter}. They do not characterize the all-inner
output in \eqref{eq:liu-envelope}.

A natural conjecture is that \eqref{eq:liu-envelope} is minimax-sharp for each
fixed $(\eta,K)$. It is false. All-inner averaging sees a fresh
without-replacement bridge in every epoch, but the epoch-end state remembers
earlier bridges. Before curvature has time to rectify these fluctuations,
the two effects cancel more strongly than the envelope records. The correct
stochastic law for $\alpha=\eta nL\leq1$ is
\begin{equation}
 \eta^2nL\sigma_\star^2
 \left[\frac1K+\min\{(\alpha K)^2,1\}\right].       \label{eq:small-law}
\end{equation}
It has three phases:
\begin{center}
\begin{tabular}{@{}lll@{}}
\toprule
horizon & dominant mechanism & stochastic error \\
\midrule
$K\lesssim\alpha^{-2/3}$ & fresh bridge
  & $\eta^2nL\sigma_\star^2/K$ \\
$\alpha^{-2/3}\lesssim K\lesssim\alpha^{-1}$ & rectification
  & $\eta^2nL\sigma_\star^2(\alpha K)^2$ \\
$K\gtrsim\alpha^{-1}$ & rectified floor
  & $\eta^2nL\sigma_\star^2$ \\
\bottomrule
\end{tabular}
\end{center}
Thus more epochs can temporarily make the fixed-step minimax error worse.
The minimum is smaller than the eventual floor by $\alpha^{2/3}$.

Our proof isolates the two mechanisms. For the upper bound, freeze every
component gradient at a minimizer. The resulting permutation prefixes reset
to zero after each epoch and have an exact bridge-area variance of
$\eta^2(n+1)\sigma_\star^2/12$. Every nonlinear and state-memory effect is a
single Lipschitz remainder. A prefix-sum operator controls this remainder
when $\alpha K$ is small, and the general bound of \citet{liu2026dominates}
covers the remaining horizons.

For the lower bound, a centered common-Hessian quadratic generates the
$1/K$ bridge term. It also satisfies an exact identity that makes its
all-inner average proportional to its final state, so its error eventually
decays rather than reaching a floor. A piecewise quadratic adapted from the
one-epoch construction of \citet{cha2023tighter} rectifies the bridge into a
mean displacement of order
$\eta\sigma_\star\sqrt n\min\{\alpha K,1\}$. Choosing the larger of these
two gadgets and adding one deterministic coordinate gives a matching
two-dimensional instance. Curvature rescaling handles large nominal steps
without padding or inactive-component arguments.

\paragraph{Contributions.}
We give the first fixed-step minimax law for the all-inner output of random
reshuffling at every finite horizon and every stable constant step. The law
answers the lower-bound question in \citet{liu2026dominates}, but also shows
why its upper envelope is not pointwise tight. The proof introduces a reset
bridge decomposition for the all-inner transient, identifies an exact
quadratic cancellation, and transfers the known rectification gadget from
epoch endpoints to all-inner averaging.

\section{Setup}

Let
\[
 f(x)=\frac1n\sum_{i=1}^n f_i(x),
 \qquad f_i:\R^d\to\R,
\]
where every component is convex and $L$-smooth. Let $x^\star$ minimize $f$
and define
\[
 \sigma_\star^2=\frac1n\sum_{i=1}^n
       \|\nabla f_i(x^\star)\|^2.
\]
At epoch $k$, random reshuffling draws an independent uniform permutation
$\pi_k$ and performs
\begin{equation}
 x_{k,j+1}=x_{k,j}-\eta\nabla f_{\pi_k(j+1)}(x_{k,j}),
 \quad j=0,\ldots,n-1,
 \quad x_{k+1,0}=x_{k,n}.                 \label{eq:rr}
\end{equation}
We study the pre-update all-inner average
\begin{equation}
 \bar x_K=\frac1{nK}\sum_{k=0}^{K-1}\sum_{j=0}^{n-1}x_{k,j}.       \label{eq:average}
\end{equation}
For a constant step, this is exactly the weighted output in Theorem~1 of
\citet{liu2026dominates}.

For $D,\sigma\geq0$, define the fixed-step minimax risk
\begin{equation}
 \Risk_{n,K}(\eta,L,D,\sigma)
 =\sup \E[f(\bar x_K)-f(x^\star)],                       \label{eq:risk}
\end{equation}
where the supremum ranges over all dimensions, all finite sums satisfying the
assumptions above, and all initial points with
$\|x_{0,0}-x^\star\|\leq D$ and $\sigma_\star\leq\sigma$.
The lower-bound instance may depend on $(\eta,K)$, as usual in a minimax
statement.

Set
\begin{equation}
 \alpha=\eta nL,
 \qquad \delta=\min\{\alpha,1\},
 \qquad
 \Stoch_{n,K}=\eta\sigma^2\delta
 \left[\frac1K+\min\{(\delta K)^2,1\}\right].            \label{eq:S}
\end{equation}
For $\alpha\geq1$, the bracket lies in $[1,2]$, so
$\Stoch_{n,K}\asymp\eta\sigma^2$.

\section{The finite-horizon minimax law}

\begin{theorem}[All steps and all horizons]\label{thm:main}
There are universal constants $c,C>0$ such that for every $n\geq8$,
$K\geq1$, $L>0$, $D,\sigma\geq0$, and $0<\eta\leq1/(6L)$,
\begin{align}
 c\left[
  \min\left\{LD^2,\frac{D^2}{\eta nK}\right\}
  +\Stoch_{n,K}\right]
 &\leq \Risk_{n,K}(\eta,L,D,\sigma) \nonumber\\
 &\leq C\left[
  \min\left\{LD^2,\frac{D^2}{\eta nK}\right\}
  +\Stoch_{n,K}\right].                         \label{eq:main}
\end{align}
The lower bound is attained in dimension at most two. For even $n$ it uses
balanced signs. For odd $n$ the same construction uses one zero sign.
\end{theorem}

The theorem separates a deterministic optimization term from a stochastic
reshuffling term. It does not assert that every instance follows the
three-phase law. In fact, the exact quadratic identity in
Section~\ref{sec:lower} gives an instance whose error vanishes after the
bridge phase. The statement is a worst-case characterization at each fixed
step and horizon.

The result also explains why the tuned rate in \citet{liu2026dominates} is
correct despite the slack in \eqref{eq:liu-envelope}. Minimizing
\eqref{eq:main} over stable constant steps gives, up to universal constants,
\begin{equation}
 \frac{LD^2}{nK}
 +\min\left\{
   \frac{\sigma D}{\sqrt{nK}},
   \left(\frac{L\sigma^2D^4}{nK^2}\right)^{1/3}
 \right\}.                                             \label{eq:tuned}
\end{equation}
The transient dip concerns fixed steps. Tuning can move the trajectory to a
different phase.

\section{Upper bound: a reset bridge and one remainder}
\label{sec:upper}

Translate $x^\star$ to zero and write
\[
 g_i=\nabla f_i(0),
 \qquad h_i(x)=\nabla f_i(x)-g_i.
\]
Then $\sum_i g_i=0$ and $\|h_i(x)\|\leq L\|x\|$. First run
\eqref{eq:rr} from zero and call this trajectory $z$. Flatten its $N=nK$
pre-update iterates into $z_0,\ldots,z_{N-1}$. Inside each epoch define the
reset bridge
\begin{equation}
 Y_{k,j}=-\eta\sum_{r=1}^j g_{\pi_k(r)},
 \qquad j=0,\ldots,n-1.                              \label{eq:bridge}
\end{equation}
Completed epochs contribute zero to the frozen gradients. Therefore, after
flattening $Y_{k,j}$ into $Y_t$,
\begin{equation}
 z_t=Y_t+R_t,
 \qquad
 R_t=-\eta\sum_{s<t}h_{i_s}(z_s).                    \label{eq:remainder}
\end{equation}

The bridge has exact finite-population moments.
\begin{lemma}[Reset bridge]\label{lem:bridge}
For $j=0,\ldots,n-1$,
\begin{align}
 \E\|Y_{k,j}\|^2
 &=\eta^2\sigma_\star^2\frac{j(n-j)}{n-1},            \label{eq:prefix-moment}\\
 \E\left\|\frac1N\sum_{t=0}^{N-1}Y_t\right\|^2
 &=\frac{\eta^2(n+1)}{12K}\sigma_\star^2.            \label{eq:area-moment}
\end{align}
Also,
\begin{equation}
 \frac1N\sum_{t=0}^{N-1}\E\|Y_t\|^2
 =\frac{\eta^2(n+1)}6\sigma_\star^2.                 \label{eq:path-moment}
\end{equation}
\end{lemma}

The $1/K$ term in \eqref{eq:S} is exactly the variance reduction from
averaging the independent centered bridge areas in
\eqref{eq:area-moment}. The second term comes from controlling $R$.

\begin{lemma}[Short-horizon memory]\label{lem:memory}
Let $\theta=\eta LN=\alpha K$. If $\theta\leq1/2$, then
\begin{equation}
 \E\|\bar z_K\|^2
 \leq C_0\eta^2n\sigma_\star^2
 \left(\frac1K+\theta^2\right)                       \label{eq:short-distance}
\end{equation}
for a numerical constant $C_0$.
\end{lemma}

The proof is a one-line operator estimate. Let $P$ be the strictly lower
triangular prefix-sum matrix on $N$ coordinates. With
$a_t=\|z_t\|$ and $b_t=\|Y_t\|$, \eqref{eq:remainder} gives
$a\leq b+\eta LPa$. Since $\|P\|_{2\to2}\leq N$,
\[
 \|a\|_2\leq2\|b\|_2,
 \qquad
 \|R\|_2\leq2\theta\|b\|_2.
\]
Jensen's inequality, \eqref{eq:area-moment}, and
\eqref{eq:path-moment} prove \eqref{eq:short-distance}. This argument uses
no Hessians, commutativity, or dimension restriction.

Smoothness of $f$ now gives
\begin{equation}
 \E[f(\bar z_K)-f^\star]
 \leq C\eta^2nL\sigma_\star^2
 \left(\frac1K+(\alpha K)^2\right),
 \qquad \alpha K\leq\tfrac12.                        \label{eq:short-function}
\end{equation}
To restore an arbitrary initialization, couple $x$ and $z$ with the same
permutations. Every component gradient map is nonexpansive when
$\eta\leq2/L$, so $\|\bar x_K-\bar z_K\|\leq D$. Co-coercivity and
smoothness imply
\begin{equation}
 f(\bar x_K)-f^\star
 \leq2\bigl(f(\bar z_K)-f^\star\bigr)+LD^2.           \label{eq:couple}
\end{equation}
When $\alpha K\leq1/2$, the deterministic minimum in
\eqref{eq:main} equals $LD^2$, so this proves the desired short-horizon
bound.

For $\alpha K>1/2$, Theorem~1 of \citet{liu2026dominates} gives
\begin{equation}
 \E[f(\bar x_K)-f^\star]
 \leq \frac{6D^2}{\eta nK}
 +51\eta\sigma_\star^2\min\{\alpha,1\}.               \label{eq:liu}
\end{equation}
If $\alpha\leq1$, the bracket in \eqref{eq:S} is at least $1/4$ in this
region. If $\alpha\geq1$, it is at least one. The first term in
\eqref{eq:liu} is within a universal factor of the capped deterministic term.
This completes the upper bound.

\section{Lower bound: cancellation and rectification}
\label{sec:lower}

The deterministic term follows from identical components
$f_i(w)=\lambda w^2/2$, initialized at $w_0=D$, with
$\lambda=\min\{L,(\eta nK)^{-1}\}$. Their geometric all-inner average is a
constant fraction of $D$, which yields
\begin{equation}
 f(\bar w_K)-f^\star
 \geq c\min\left\{LD^2,\frac{D^2}{\eta nK}\right\}.   \label{eq:det-lower}
\end{equation}

It remains to obtain \eqref{eq:S}. We first expose the fresh bridge with a
centered common-Hessian quadratic. Let $s_i$ be balanced signs and set
\begin{equation}
 Q_i(u)=\frac\ell2u^2+a s_i u,
 \qquad u_0=0.                                        \label{eq:centered}
\end{equation}
For odd $n$, one sign is zero. The scale $a$ is adjusted so that
$n^{-1}\sum_i a^2s_i^2=\sigma^2$. Put $q=1-\eta\ell$ and $r=q^n$. Summing
all $nK$ update equations and using $\sum_i s_i=0$ gives the pathwise identity
\begin{equation}
 \boxed{\quad \bar u_K=-\frac{u_{K,0}}{\eta\ell nK}.\quad}             \label{eq:telescoping}
\end{equation}
This exact cancellation is special to all-inner averaging.

If
\[
 V_n(q)=\sum_{j=1}^n\left(q^{n-j}-\frac1n
             \sum_{r=1}^nq^{n-r}\right)^2,
\]
finite-population covariance gives
\begin{equation}
 \E[Q(\bar u_K)-Q(0)]
 =\frac{\sigma^2V_n(q)}{2\ell n(n-1)K^2}
   \frac{1-r^{2K}}{1-r^2}                              \label{eq:exact-quadratic}
\end{equation}
for every $n$, where $Q=n^{-1}\sum_iQ_i=\ell u^2/2$.
When $\beta=\eta n\ell$ and $\beta K$ are small,
$V_n(q)\asymp(\eta\ell)^2n^3$ and the geometric factor is
$\asymp K$. Hence
\begin{equation}
 \E[Q(\bar u_K)-Q(0)]
 \geq c\frac{\eta^2n\ell\sigma^2}{K}.                \label{eq:bridge-lower}
\end{equation}
For $K\gg1/\beta$, \eqref{eq:exact-quadratic} instead decays as
$K^{-2}$. Common-Hessian quadratics cannot create the floor.

The floor comes from curvature rectification. Consider
\begin{equation}
 r_i(v)=\frac12\left(
   \ell\1_{\{v<0\}}+\frac\ell{2415}\1_{\{v\geq0\}}
 \right)v^2+a s_i v,
 \qquad v_0=0.                                        \label{eq:rectifier}
\end{equation}
Every component is convex and $\ell$-smooth. The average objective is
minimized at zero and is $\ell/2415$-strongly convex.

\begin{lemma}[Finite-horizon rectification]\label{lem:rectifier}
There are universal constants $b,c_1,c_2>0$ such that, whenever
$\beta=\eta n\ell\leq b$, random reshuffling on
\eqref{eq:rectifier} satisfies, with
$B=\eta a\sqrt n$,
\begin{equation}
 m_{k+1}\geq(1-c_1\beta)m_k+c_2\beta B,
 \qquad m_k=\E[v_{k,0}],\quad m_0=0.                  \label{eq:recurrence}
\end{equation}
Consequently, if $K\geq\beta^{-2/3}$, then
\begin{equation}
 \E[r(\bar v_K)-r(0)]
 \geq c\eta^2n\ell a^2\min\{(\beta K)^2,1\}.        \label{eq:rect-lower}
\end{equation}
The statement holds for balanced signs when $n$ is even and for balanced
signs plus one zero when $n$ is odd.
\end{lemma}

The recurrence is the one-epoch rectification estimate proved in
Appendix~B of \citet{cha2023tighter}. The new all-inner transfer is short.
Couple \eqref{eq:rectifier} to the common-curvature $\ell$ quadratic under
the same permutation. Monotonicity gives a pathwise lower comparison.
Conditioned on an epoch start,
\[
 \E[v_{k,j}\mid v_{k,0}]\geq(1-\eta\ell)^jv_{k,0}.
\]
Solving \eqref{eq:recurrence} gives
$m_k\geq cB\min\{\beta k,1\}$. Averaging over the starts of the first $K$
epochs, and using $K\geq\beta^{-2/3}>1$, yields
$\E\bar v_K\geq cB\min\{\beta K,1\}$. Strong convexity and Jensen's
inequality prove \eqref{eq:rect-lower}. Appendix~\ref{app:rectifier}
records the constants and the odd-$n$ reduction.

To cover every nominal step, choose an effective curvature
\begin{equation}
 \ell=\frac{\beta}{\eta n},
 \qquad \beta=\min\{\alpha,b\}.                       \label{eq:rescale}
\end{equation}
Then $\ell\leq L$. Use \eqref{eq:centered} when
$K\lesssim\beta^{-2/3}$ and \eqref{eq:rectifier} otherwise. The two bounds
give
\[
 c\eta\sigma^2\beta
 \left[\frac1K+\min\{(\beta K)^2,1\}\right].
\]
Since $\beta/\delta\geq b$ and, for $0<\beta\leq\delta\leq1$,
\[
 \frac1K+\min\{(\beta K)^2,1\}
 \geq\left(\frac\beta\delta\right)^2
 \left[\frac1K+\min\{(\delta K)^2,1\}\right],
\]
this is a universal constant fraction of \eqref{eq:S}. Finally, take the
orthogonal direct sum of the deterministic coordinate and the selected
stochastic coordinate. Block-diagonal smoothness remains at most $L$, the
initial distance is exactly $D$, and the gradient variance at the minimizer is
exactly $\sigma^2$. This proves the lower bound in Theorem~\ref{thm:main}.

\section{Relation to prior reshuffling theory}

Classical analyses established improved rates for random reshuffling under
small-step, large-epoch, strong-convexity, or bounded-gradient conditions
\citep{gurbuzbalaban2021why,nagaraj2019sgd,haochen2019random,
mishchenko2020random,nguyen2021unified}. Lower bounds for final iterates and
quadratic problems reveal condition-number and epoch transitions
\citep{safran2020how,safran2021many}. Last-iterate and primal-dual analyses
have since sharpened the output guarantees
\citep{liu2024last,cai2024tighter}.

The closest lower bound is \citet{cha2023tighter}. It treats arbitrary
nonnegative combinations of epoch-boundary iterates. Its convex corollary
uses a small step and a sufficiently large horizon. We use its one-epoch
piecewise-quadratic estimate, but our output contains every inner iterate and
our result tracks every finite horizon. The bridge term and the transient
dip have no analogue in an endpoint-only average.

The closest upper bound is \citet{liu2026dominates}. It introduced the exact
all-inner output in \eqref{eq:average}, proved \eqref{eq:liu} at all stable
steps and horizons, and explicitly left the corresponding complete lower
characterization open. Our theorem closes that direction. It also shows
that the stochastic part of \eqref{eq:liu} is loose by an unbounded factor at
$K\asymp\alpha^{-2/3}$ as $\alpha\to0$.

Constant-step stationary-distribution and bias expansions study a different
asymptotic object, usually epoch endpoints under strong monotonicity or
additional structure \citep{emmanouilidis2026shuffling}. They are
complementary to the finite-horizon convex minimax law here.

\section{Discussion}

The nonmonotonicity is an output effect, not a claim that running a fixed
instance longer must hurt. Minimax instances can depend on the horizon. The
law says which mechanism an adversary can activate at each horizon. It also
clarifies the role of state memory. Memory creates no fourth minimax scale.
During the transient it is contained by the Lipschitz remainder, while in the
common-Hessian lower gadget it cancels the bridge exactly.

Our theorem uses the maximum component smoothness $L$. A sharper law that
separates average and maximum smoothness, as in the nonuniform upper theorem
of \citet{liu2026dominates}, remains open. The result also concerns the
specific all-inner average. Other output rules can have different horizon
laws.

\section*{Reproducibility statement}
All assumptions and complete proofs are included in the manuscript. The
supplementary verification script checks the exact quadratic identity by
permutation enumeration and Monte Carlo simulation, including the pre-update
indexing convention. No external data or specialized hardware is required.

\section*{AI use statement}
Generative AI tools were used for exploratory theorem generation, literature
search, adversarial proof auditing, numerical identity checks, and language
editing. The authors rederived every displayed result, read the cited primary
papers, ran the verification artifact, and checked the stated assumptions and
boundary cases. The authors take full responsibility for the claims and the
final manuscript.

\bibliography{references}
\bibliographystyle{iclr2027_conference}

\appendix

\section{Finite-population identities}\label{app:bridge}

We prove Lemma~\ref{lem:bridge}. Let $G_1,\ldots,G_n\in\R^d$ satisfy
$\sum_iG_i=0$ and $n^{-1}\sum_i\|G_i\|^2=\sigma_G^2$. If $\pi$ is uniform,
then
\[
 \E\langle G_{\pi(r)},G_{\pi(s)}\rangle
 =\begin{cases}
 \sigma_G^2,&r=s,\\
 -\sigma_G^2/(n-1),&r\neq s.
 \end{cases}
\]
It follows directly that
\begin{equation}
 \E\left\|\sum_{r=1}^jG_{\pi(r)}\right\|^2
 =\sigma_G^2\frac{j(n-j)}{n-1}.                       \label{eq:finite-prefix}
\end{equation}
Taking $G_i=g_i$ proves \eqref{eq:prefix-moment}. Summing
\eqref{eq:finite-prefix} over $j=0,\ldots,n-1$ and using
\[
 \sum_{j=0}^{n-1}j(n-j)=\frac{n(n-1)(n+1)}6
\]
proves \eqref{eq:path-moment}.

For the area of one epoch,
\[
 A_\pi=\frac1n\sum_{j=0}^{n-1}Y_j
 =-\frac\eta n\sum_{r=1}^{n-1}(n-r)g_{\pi(r)}.
\]
For arbitrary scalar weights $w_r$, the covariance above gives
\begin{equation}
 \E\left\|\sum_{r=1}^nw_rg_{\pi(r)}\right\|^2
 =\frac{n\sigma_\star^2}{n-1}
   \sum_{r=1}^n(w_r-\bar w)^2.                         \label{eq:weighted-fp}
\end{equation}
With $w_r=n-r$ and
$\sum_r(w_r-\bar w)^2=n(n^2-1)/12$, we obtain
\[
 \E\|A_\pi\|^2=\frac{\eta^2(n+1)}{12}\sigma_\star^2.
\]
The epoch areas are independent and centered. Their average has variance
$1/K$ times this quantity, proving \eqref{eq:area-moment}.

\section{Complete upper-bound proof}\label{app:upper}

We first prove Lemma~\ref{lem:memory}. Index the $N=nK$ pre-update iterates
from zero. Let $P\in\R^{N\times N}$ have entries $P_{ts}=1$ if $s<t$ and
zero otherwise. Equation~\eqref{eq:remainder} and
$\|h_i(z)\|\leq L\|z\|$ imply componentwise
\[
 a\leq b+\eta LPa,
 \qquad a_t=\|z_t\|,
 \quad b_t=\|Y_t\|.
\]
Since $\|P\|_{2\to2}\leq\|P\|_{\mathrm F}<N$ and
$\theta=\eta LN\leq1/2$,
\begin{equation}
 \|a\|_2\leq\|b\|_2+\theta\|a\|_2
 \leq2\|b\|_2.                                      \label{eq:a-b}
\end{equation}
Likewise, the sequence of remainder norms satisfies
\begin{equation}
 \|R\|_2\leq\eta L\|Pa\|_2
 \leq2\theta\|b\|_2.                                \label{eq:R-b}
\end{equation}
Jensen's inequality gives
\[
 \E\|\bar R\|^2
 \leq\frac1N\E\|R\|_2^2
 \leq4\theta^2\frac1N\E\|b\|_2^2
 =\frac{2(n+1)}{3}\eta^2\theta^2\sigma_\star^2.
\]
Together with \eqref{eq:area-moment} and
$\|u+v\|^2\leq2\|u\|^2+2\|v\|^2$, this proves
\eqref{eq:short-distance} with, for example, $C_0=4$ for $n\geq2$.

For completeness, we justify the initialization coupling. If $T_i(x)=
x-\eta\nabla f_i(x)$, co-coercivity gives
\begin{align*}
 \|T_i(x)-T_i(y)\|^2
 &\leq\|x-y\|^2
 -(2\eta/L-\eta^2)
   \|\nabla f_i(x)-\nabla f_i(y)\|^2.
\end{align*}
Thus $T_i$ is nonexpansive for $\eta\leq2/L$. Coupled trajectories remain
within $D$, and so do their averages. For any $u,v$, smoothness gives
\[
 f(u)-f^\star
 \leq f(v)-f^\star+\langle\nabla f(v),u-v\rangle
       +\frac L2\|u-v\|^2.
\]
The standard inequality $\|\nabla f(v)\|^2\leq2L(f(v)-f^\star)$ and Young's
inequality yield \eqref{eq:couple}.

If $\theta\leq1/2$, then
$D^2/(\eta nK)=LD^2/\theta\geq2LD^2$, so the capped deterministic term is
$LD^2$. Equations \eqref{eq:short-function} and \eqref{eq:couple} give the
upper half of \eqref{eq:main} in this region.

If $\theta>1/2$, apply \eqref{eq:liu}. When $\theta\in(1/2,1)$,
$D^2/(\eta nK)<2LD^2$, and when $\theta\geq1$ it is the smaller deterministic
term. Hence the first term in \eqref{eq:liu} is at most twelve times the
capped term. If $\alpha\leq1$, then
\[
 K^{-1}+\min\{(\alpha K)^2,1\}\geq1/4.
\]
If $\alpha\geq1$, then $\delta=1$ and the same bracket is at least one. This
absorbs the stochastic term in \eqref{eq:liu} and completes the proof.

\section{Exact common-Hessian cancellation}\label{app:quadratic}

We derive \eqref{eq:exact-quadratic} and its lower consequence. First assume
$n$ is even and $s_i\in\{-1,1\}$ are balanced. For one epoch,
\[
 u_{k+1,0}=r u_{k,0}+\xi_k,
 \qquad
 \xi_k=-\eta a\sum_{j=1}^nq^{n-j}s_{\pi_k(j)}.
\]
The innovations are independent and centered. Equation
\eqref{eq:weighted-fp} gives
\[
 \E\xi_k^2=\frac{\eta^2na^2}{n-1}V_n(q).
\]
Starting from zero,
\[
 \E u_{K,0}^2
 =\frac{\eta^2na^2}{n-1}V_n(q)
   \sum_{k=0}^{K-1}r^{2k}.
\]
The sum of the $nK$ scalar update equations is
\[
 u_{K,0}-u_{0,0}
 =-\eta\ell\sum_{k,j}u_{k,j}
   -\eta a\sum_{k,j}s_{\pi_k(j)}.
\]
The last term vanishes epoch by epoch. This proves
\eqref{eq:telescoping}. Substitution into
$Q(\bar u_K)-Q(0)=\ell\bar u_K^2/2$ proves
\eqref{eq:exact-quadratic}.

We record explicit bounds. If $\beta=\eta n\ell\leq1/4$, then
$q^{n-1}\geq1-\beta\geq3/4$. The pairwise variance identity gives
\begin{align*}
 V_n(q)
 &=\frac1{2n}\sum_{p,r=0}^{n-1}(q^p-q^r)^2\\
 &\geq\frac{q^{2n-2}(\eta\ell)^2}{2n}
        \sum_{p,r=0}^{n-1}(p-r)^2
 \geq c(\eta\ell)^2n^3.                              \label{eq:V-lower}
\end{align*}
If also $\beta K\leq1/4$, then $r^{2k}\geq(1-\beta)^{2K}\geq c$ for
$k<K$, so $\sum_{k<K}r^{2k}\geq cK$. Equation
\eqref{eq:exact-quadratic} now proves \eqref{eq:bridge-lower}.

For odd $n$, choose $(n-1)/2$ signs of each type and one zero. Put
$\rho_n=n^{-1}\sum_i s_i^2=(n-1)/n$ and $a=\sigma/\sqrt{\rho_n}$.
The weighted covariance formula has $a^2\rho_n=\sigma^2$ in place of $a^2$,
so the same calculation and constants apply.

\section{The rectification lemma}\label{app:rectifier}

We give the promised reduction to the proved one-epoch estimates of
\citet{cha2023tighter}. Their Appendix~B considers exactly
\eqref{eq:rectifier} for even $n$, with $a$ denoted by $\nu$. Lemmas B.2--B.4
and the calculation following them show, for
$\beta\leq1/161$, that an epoch starting from a fixed $v\geq0$ satisfies
\begin{equation}
 \E[v^+\mid v]
 \geq(1-3\beta/4)v+\frac{\eta^2\ell a n^{3/2}}{18000},
                                                               \label{eq:cha-pos}
\end{equation}
while a negative start satisfies
\begin{equation}
 \E[v^+\mid v]\geq(1-160\beta/161)v.                  \label{eq:cha-neg}
\end{equation}
They also prove that, from a nonnegative initialization, every epoch start is
nonnegative with probability at least $1/2$. Since $v<0$ reverses the
comparison between the two contraction coefficients, conditioning on the
sign of the start gives
\begin{equation}
 m_{k+1}\geq(1-3\beta/4)m_k+\frac{\beta B}{36000},
 \qquad B=\eta a\sqrt n.                              \label{eq:m-exact}
\end{equation}
This is \eqref{eq:recurrence}. Taking a smaller universal $b$ if needed,
iteration from $m_0=0$ yields
\begin{equation}
 m_k\geq cB\min\{\beta k,1\}.                         \label{eq:m-solve}
\end{equation}

We next transfer endpoint drift to the all-inner output. Let $\widetilde v_j$
be the trajectory with the same linear signs and common curvature $\ell$.
For $\eta\ell\leq1$, both component update maps are nondecreasing and
\[
 T_i^{\mathrm{rect}}(x)\geq T_i^{\mathrm{quad}}(x)
 \quad\text{for every }x.
\]
Induction gives $v_j\geq\widetilde v_j$ pathwise. Conditional on the epoch
start, uniform permutation marginals and balanced signs give
\[
 \E[\widetilde v_j\mid v_0]=(1-\eta\ell)^jv_0.
\]
Therefore
\[
 \E v_{k,j}\geq(1-\eta\ell)^jm_k\geq(1-\beta)m_k.
\]
If $K\geq\beta^{-2/3}$ and $b$ is small, then $K\geq2$. Averaging
\eqref{eq:m-solve} over $k=0,\ldots,K-1$ gives
\begin{equation}
 \E\bar v_K\geq cB\min\{\beta K,1\}.                 \label{eq:mean-all}
\end{equation}
The average objective in \eqref{eq:rectifier} is
$\ell/2415$-strongly convex and minimized at zero. Jensen's inequality gives
\[
 \E[r(\bar v_K)-r(0)]
 \geq\frac\ell{4830}(\E\bar v_K)^2,
\]
which proves \eqref{eq:rect-lower}.

For odd $n$, use balanced $+1$ and $-1$ signs and one zero. We explain why
the cited one-epoch proof still applies. If $S_j$ is a prefix sum of a
uniform permutation of these signs, sampling without replacement gives, for
$j\leq n/2$,
\[
 \E S_j^2\geq c j,
 \qquad \E S_j^4\leq Cj^2.
\]
H\"older's inequality and Paley--Zygmund imply
$\E|S_j|\geq c\sqrt j$ and $\Pr(S_j>0)=\Pr(S_j<0)\geq c$.
These are the only sign-distribution inputs in Lemmas B.2--B.4 of
\citet{cha2023tighter}. Repeating their displayed inequalities with the new
universal constants proves \eqref{eq:cha-pos}--\eqref{eq:cha-neg} up to
constant factors. Scaling $a$ by $\sqrt{n/(n-1)}$ preserves
$n^{-1}\sum_i a^2s_i^2=\sigma^2$ exactly.

\section{Completing the lower bound}\label{app:lower}

Choose a universal $b>0$ small enough that $b\leq1/161$ and
$b^{1/3}\leq1/4$. Let $\beta=\min\{\alpha,b\}$ and choose $\ell$ as in
\eqref{eq:rescale}. Then $\ell\leq L$ and $\eta n\ell=\beta$.

If $K\leq\beta^{-2/3}$, then $\beta K\leq\beta^{1/3}\leq1/4$ and
Appendix~\ref{app:quadratic} gives
\[
 \E[Q(\bar u_K)-Q(0)]
 \geq c\eta^2n\ell\sigma^2/K.
\]
In this regime, $(\beta K)^2\leq1/K$, so this is a constant fraction of
\[
 \eta^2n\ell\sigma^2
 \left[K^{-1}+\min\{(\beta K)^2,1\}\right].
\]
If $K\geq\beta^{-2/3}$, Lemma~\ref{lem:rectifier} applies and
$K^{-1}\leq(\beta K)^2$. Its lower bound is again a constant fraction of
the same display. Thus one of the two one-dimensional stochastic gadgets
always supplies
\begin{equation}
 c\eta\sigma^2\beta
 \left[K^{-1}+\min\{(\beta K)^2,1\}\right].           \label{eq:beta-law}
\end{equation}

Let $\delta=\min\{\alpha,1\}$. We have
$\beta/\delta\geq b$. For $r=\beta/\delta\leq1$,
\[
 \min\{(\beta K)^2,1\}
 \geq r^2\min\{(\delta K)^2,1\},
 \qquad K^{-1}\geq r^2K^{-1}.
\]
Hence \eqref{eq:beta-law} is at least $cb^3\Stoch_{n,K}$.

Finally, verify \eqref{eq:det-lower}. Put $T=nK$ and
$\lambda=\min\{L,(\eta T)^{-1}\}$. The pre-update iterates are
$w_t=(1-\eta\lambda)^tD$, and
\[
 \bar w_K
 =D\frac{1-(1-\eta\lambda)^T}{T\eta\lambda}.
\]
If $\eta LT\leq1$, then $\lambda=L$ and $\bar w_K\geq cD$. If
$\eta LT\geq1$, then $\lambda=(\eta T)^{-1}$ and again
$\bar w_K\geq cD$. Multiplying by $\lambda/2$ proves
\eqref{eq:det-lower}. The orthogonal direct sum described in the main text
finishes the proof of Theorem~\ref{thm:main}.

\section{Optimizing the step}\label{app:tuning}

We briefly derive \eqref{eq:tuned}. A stable step of order $1/L$ contributes
the baseline $LD^2/(nK)$. In the regime $\alpha\geq1$, balancing
$D^2/(\eta nK)$ with $\eta\sigma^2$ gives
$\sigma D/\sqrt{nK}$. In the small-step regime after rectification has
saturated, balancing $D^2/(\eta nK)$ with
$\eta^2nL\sigma^2$ gives
$(L\sigma^2D^4/(nK^2))^{1/3}$. Taking the better admissible balance and the
stability boundary yields \eqref{eq:tuned}. The lower half of
Theorem~\ref{thm:main} shows that the same expression is minimax-sharp for
constant-step random reshuffling.

\end{document}
